\documentclass[12pt]{article}
\usepackage[english]{babel}
\usepackage[T1]{fontenc}
\usepackage{carlito}  % Calibri-like font
\usepackage{newtxtext,newtxmath}  % Times New Roman-like font
\usepackage{graphicx}
\usepackage{setspace}
\usepackage{cite}
\usepackage{amsmath}
\usepackage{amsfonts}
% \usepackage{amssymb}
% \let\Bbbk\relax
\usepackage{newtxtext,newtxmath}
\usepackage{algorithm}
\usepackage{algorithmic}
\usepackage{graphicx}
\usepackage{caption}
\usepackage{subcaption}
\usepackage{hyperref}
\usepackage{booktabs}
\usepackage{multirow}
\usepackage{enumitem}
\usepackage{adjustbox}
\usepackage{array}
\usepackage{titlesec}
\usepackage[left=2.5cm,right=2.5cm,top=2.5cm,bottom=2.5cm]{geometry}
\usepackage{indentfirst}
\usepackage{lipsum}

% \newcommand{\Bbk}

% Mengatur indentasi paragraf
\setlength{\parindent}{0.6cm}

% Menghapus jarak antar paragraf
\setlength{\parskip}{0pt}

% Mengatur format judul section dengan Times New Roman-like font
\titleformat{\section}[block]{\rmfamily\bfseries}{\thesection}{1em}{}

% Menghapus jarak sebelum dan sesudah section
\titlespacing*{\section}{0pt}{0pt}{0pt}

% Command untuk mengubah ke Times New Roman-like font
\newcommand{\headingfont}{\rmfamily}

% Mengatur font default ke Calibri-like
\renewcommand{\familydefault}{\sfdefault}

\begin{document}

\begin{center}
{\headingfont\fontsize{18}{22}\selectfont\textbf{Journal of Science and Informatics for Society (JSIS)}}
\end{center}

\begin{center}
\headingfont\textbf{Author}

instances\\
Email
\end{center}

\begin{center}
\headingfont\textbf{\textit{Abstract}}
\end{center}

\doublespacing
\noindent
\textit{Lorem ipsum dolor sit amet, consectetur adipiscing elit, sed do eiusmod tempor incididunt ut labore et dolore magna aliqua. Ut enim ad minim veniam, quis nostrud exercitation ullamco laboris nisi ut aliquip ex ea commodo consequat. Duis aute irure dolor in reprehenderit in voluptate velit esse cillum dolore eu fugiat nulla pariatur. Excepteur sint occaecat cupidatat non proident, sunt in culpa qui officia deserunt mollit anim id est laborum}

\noindent

\textit{{\headingfont\textbf{Key Word:}} }

\singlespacing
\section*{\headingfont Introduction}
\lipsum[1-2]
\\

\section*{\headingfont Realted Work}
\lipsum[3-4]
\\



\bibliographystyle{IEEEtran}
\bibliography{bibtex/bib/IEEEexample}

\end{document}